\clearpage
\section{The Exact Obstruction Complex and Formal Shell Geometry}

\subsection{A joint exponent-free two-row complex}
The Smith obstruction is attached to the exponent-free phase, or shell-synthesis, rows rather than to the exponent-weighted derivative row. For every prime $p\mid abc$, define
\[
 \Lambda_p=
 \begin{cases}
 a/p,&p\mid a,\\
 b/p,&p\mid b,\\
 -c/p,&p\mid c,
 \end{cases}
 \qquad
 \mathsf w_p=
 \begin{cases}
 -ab/p,&p\mid a,\\
 ab/p,&p\mid b,\\
 0,&p\mid c.
 \end{cases}
\]
Consider
\[
 \mathsf M_{a,b,c}
 =
 \begin{pmatrix}
 \Lambda\\
 \mathsf w
 \end{pmatrix}:
 \Z^S\longrightarrow\Z^2.
\]
The first row is the exponent-free additive phase row. The second is its integer Wronskian row. Set
\[
 Q=\frac{abc}{R}.
\]

\begin{theorem}[Smith obstruction theorem]\label{thm:phase-smith}
Let $a+b=c$ be primitive, with at least two of $a,b,c$ greater than one. Then
\[
 \boxed{
 \SNF(\mathsf M_{a,b,c})
 =
 \operatorname{diag}(1,Q).
 }
\]
Consequently,
\[
 \boxed{
 \operatorname{coker}(\mathsf M_{a,b,c})_{\mathrm{tors}}
 \cong
 \Z/Q\Z.
 }
\]
The exceptional triple $(1,1,2)$ has only one nonempty prime sector, so the matrix has rank one and is excluded from the two-row statement.
\end{theorem}

\begin{proof}
The gcd of the entries of $\Lambda$ is one. Indeed, every prime dividing one sector is absent from every other nonempty sector, and an entry from another sector is not divisible by it. Hence the first Smith factor is one.

Rows restricted to one sector are proportional, so same-sector $2\times2$ minors vanish. If $p$ and $q$ lie in two different nonempty sectors, the corresponding minor is
\[
 \pm\frac{abc}{pq}
 =
 \pm Q\frac{R}{pq}.
\]
Every nonzero minor is therefore divisible by $Q$. Conversely, for each prime $\ell\mid R$, choose $p=\ell$ and choose $q$ in a different nonempty sector. Since $R$ is squarefree, $R/(pq)$ is not divisible by $\ell$. Thus the gcd of the reduced minors $R/(pq)$ is one. The gcd of all nonzero minors is exactly $Q$. The determinantal-divisor description of Smith normal form now gives the result.
\end{proof}

\subsection{The exact lattice extension}
Throughout the remaining height and exact-sequence constructions, exclude $(1,1,2)$. Define
\[
 L=\ker(\Lambda)\cap\Z^S
\]
and
\[
 L_0=L\cap\ker(\mathsf w).
\]
Because $\Lambda$ is primitive, choose $t\in\Z^S$ with $\Lambda(t)=1$. Then
\[
 \Z^S=L\oplus\Z t.
\]
The image of $\mathsf M_{a,b,c}$ is generated by $(1,\mathsf w(t))$ together with $(0,\mathsf w(L))$. Its index in $\Z^2$ is therefore the index of $\mathsf w(L)$ in $\Z$. Theorem~\ref{thm:phase-smith} shows that this index is $Q$, so
\[
 \mathsf w(L)=Q\Z.
\]
Hence
\[
 \omega(x)=\frac{\mathsf w\cdot x}{Q}
\]
is an integer-valued surjection and gives the exact sequence
\[
 \boxed{
 0\longrightarrow L_0
 \longrightarrow L
 \xrightarrow{\omega}
 \Z
 \longrightarrow0.
 }
\]
The extension splits algebraically. The difficult quantity is the norm of a section, not its existence.

For $k\in\Z$, define
\[
 H(k)
 =
 \min\left\{
 \norm{x}_\infty:
 x\in L,
 \omega(x)=k
 \right\}.
\]
Three heights must be distinguished:
\[
 H(1),
 \qquad
 \Phi=\min_{k\neq0}H(k),
 \qquad
 \overline H=
 \inf_{k\ge1}\frac{H(k)}k.
\]
The shortest nonzero detector may live in a nonprimitive obstruction coset.

\subsection{The stable real height}
Relax the integral packet to a real vector in the additive hyperplane. The exact continuous width can be computed from the three sector weights. Let
\[
 \sigma_u=\sum_{p\mid u}\frac1p
\]
and
\[
 \Delta(a,b,c)
 =
 \min\left\{
 \sigma_a+\sigma_b,
 \frac cb(\sigma_a+\sigma_c),
 \frac ca(\sigma_b+\sigma_c)
 \right\}.
\]
Then
\[
 \boxed{
 \overline H
 =
 \frac{c}{R\Delta(a,b,c)}.
 }
\]
The discrete master inequality becomes
\[
 c\le R\Delta\Phi.
\]
Since
\[
 \Phi=\overline H\,k_*\Gamma
\]
for the optimal obstruction multiple $k_*$ and an integrality factor $\Gamma\ge1$,
\[
 \boxed{
 \Delta\Phi
 =
 \frac cR k_*\Gamma.
 }
\]
This is the no-escape identity. Optimizing the same first-order invariant cannot lower the continuous floor $c/R$.

\subsection{The full shell deformation}
Introduce one variable for every valuation shell and define
\[
 \mathscr J_n(\mathbf T)
 =
 \prod_{p^e\parallel n}
 \prod_{r=0}^{e-1}(p+T_{p,r}).
\]
At $\mathbf T=0$,
\[
 \mathscr J_n(0)=n.
\]
For $a+b=c$, set
\[
 F_{a,b,c}(\mathbf T)
 =
 \mathscr J_a(\mathbf T)
 +
 \mathscr J_b(\mathbf T)
 -
 \mathscr J_c(\mathbf T).
\]
The shell gradient has entries
\[
 \frac ap,
 \qquad
 \frac bp,
 \qquad
 -\frac cp
\]
repeated according to valuation depth. The gcd of all entries is one.

\begin{theorem}[Formal affine-space theorem]
Let $d$ be the number of shell variables. After a unimodular integral change of shell coordinates, the $(\mathbf T)$-adic completion of
\[
 \Z[\mathbf T]/(F_{a,b,c})
\]
at the zero section is isomorphic to a formal power-series ring $\Z[[U_2,\ldots,U_d]]$.
\end{theorem}

\begin{proof}
A primitive gradient row can be completed to a unimodular matrix. In the new coordinates, one variable has linear coefficient one, which is a unit in $\Z$. The formal implicit-function recursion therefore solves uniquely for that variable at every degree with integral coefficients, giving the stated isomorphism of completed rings.
\end{proof}

\begin{corollary}[Integral jet extension]
Every integral first-order tangent vector to $F_{a,b,c}=0$ extends to an integral formal curve through all orders.
\end{corollary}

The consequence is decisive. Higher shell curvature exists, but there is no formal lifting obstruction to exploit. The missing arithmetic must be a global size or conductor estimate.

\subsection{Aggregate products and the integral boundary}
Over $\Q$, on the locus where all shell factors are nonzero, product-preserving directions within a sector form a torus. For $a,b>1$, the normalized products
\[
 A=\mathscr J_a/a,\qquad B=\mathscr J_b/b,\qquad C=\mathscr J_c/c
\]
satisfy $aA+bB=cC$, leaving a two-dimensional aggregate plane and internal product gauges. A unit input sector has no shell variables and leaves a one-dimensional aggregate line.

This rational description introduces denominators. It is not an assertion that the integral completion is a torus. The preceding primitive-gradient argument is the integral statement: local formal extension exists through all orders. Arithmetic size remains the separate global obligation.
